\documentclass[10pt,a4paper]{amsart}
\usepackage[margin=19mm]{geometry}
\usepackage{amsmath,amssymb,mathtools}
\usepackage[scr=rsfs,cal=euler]{mathalfa}
\usepackage{xfrac}
\usepackage[unicode,hidelinks,bookmarks=false]{hyperref}
\newcommand{\ie}{i.e.\ }
\newcommand{\cf}{cf.\ }
\newcommand{\resp}{respectively\ }
\newcommand{\Subsection}{\subsection}
\providecommand{\nobreakdash}{\nobreak\hbox{-}}
\setlength{\emergencystretch}{2em}
\multlinegap0pt
\usepackage{xcolor}
\usepackage{amssymb}
\usepackage{float}
\usepackage{tikz}
\usepackage{enumerate}
\newtheorem{proposition}{Proposition}
\newcommand{\PF}{\mathsf{PF}}
\title{Generalized parking function polytopes}
\author{Mitsuki Hanada, John Lentfer, Andr\'es R. Vindas-Mel\'endez}
\date{Source: arXiv:2212.06885v2 (CC BY 4.0)}
\begin{document}
\maketitle

\noindent\emph{Source and licence.} Generalized parking function polytopes. Version arXiv:2212.06885v2 (CC BY 4.0), distributed by arXiv under the Creative Commons Attribution 4.0 International licence, http://creativecommons.org/licenses/by/4.0/, as recorded in the arXiv licence metadata for this exact version and shown on the abstract page https://arxiv.org/abs/2212.06885v2. Full licence notice: \texttt{licenses/arxiv-2212.06885-CC-BY-4.0.txt}.\par\smallskip

\noindent\textit{Editorial scope.} Counted unit: the article's proposition proposition (source0.tex lines 634--636). The article's complete original proof of it is delivered with it (source0.tex lines 639--682, one \texttt{proof} environment of 44 lines). No mathematics is written, repaired or re-derived: the statement and the proof are the article's own lines, in the article's own notation, and no retained line is substituted or re-flowed.  Retained uncounted as the source-local context the proof needs: lines 566--568.  Nothing was omitted from any retained span, so the package carries no omission record.  No retained line contains a citation, so the wrapper carries no bibliography and no bibliography entry.  No retained line contains a figure environment, an image-inclusion command or drawing code, so no image is delivered and the package declares no asset dependency.  Editorial formatting only, in addition: an AMS \texttt{amsart} preamble that restates the article's own declarations and macros used by the retained lines, namely the environments \texttt{proposition} on separate counters, together with the standard packages those lines need.  The retained environments print in this document as Proposition 1 in order of appearance, whereas the article prints the same items with the numbering of its own full document.  The complete native source of the article as distributed in the arXiv e-print for this exact version is bundled unchanged as \texttt{sources/134849/source0.tex}.
\begin{proposition}\label{prop:permutahedron}
The regular permutahedron appears as a facet of the classical parking function polytope exactly once. 
\end{proposition}

\begin{proposition}
The $(n-1)$-dimensional parking function polytope $\PF_{n-1}$ appears as a facet of the $n$-dimensional parking function polytope $\PF_n$ exactly $n$ times.
\end{proposition}

\begin{proof}
Consider the inequality description of $\PF_n$:
\begin{align}
    1 \leq x_i &\leq n, &\text{ for } 1 \leq i \leq n,\label{eq:1}\\
    x_i+x_j &\leq n + (n-1), &\text{ for } i < j,\label{eq:2}\\
    x_i+x_j+x_k &\leq n + (n-1) + (n-2), &\text{ for } i < j < k,\label{eq:3}\\
    &\vdots\nonumber\\
    x_{i_1}+ x_{i_2}+ \cdots + x_{i_{n-2}} &\leq n + (n-1) + \cdots + 3, &\text{ for } i_1 < i_2 < \cdots < i_{n-2},\label{eq:n-2}\\
    x_{1}+ x_{2}+ \cdots + x_{n} &\leq n + (n-1) + \cdots + 1.\label{eq:n}
\end{align}
Fix $x_i = n$ for some $i$; without loss of generality say $i=n$. 
Then we proceed to reduce the system of inequalities.
For all $1 \leq i \leq n-1$, we still have $1 \leq x_i$, but if any \[x_i > n - 1, \text{ then } x_i + x_n >  n + (n-1),\] which contradicts (\ref{eq:2}).
Thus, we have $1 \leq x_i \leq n-1$ for $1 \leq i \leq n-1$. 
Next, for $i < j < n$, if \[x_i + x_j > (n-1) + (n-2), \text{ then } x_i + x_j + x_n > n + (n-1) + (n-2),\] which contradicts (\ref{eq:3}). 
Thus, we have $x_i + x_j \leq (n-1) + (n-2)$ for $i < j < n$. 
Continuing this process, we can refine the inequalities up through: for $i_1 < i_2 < \cdots < i_{n-3} < n$, if \[x_{i_1} + x_{i_2} + \cdots + x_{i_{n-3}} > (n-1) + (n-2) + \cdots + 3,\] then \[x_{i_1} + x_{i_2} + \cdots + x_{i_{n-3}} + x_n > n + (n-1) + (n-2) + \cdots +3,\] which contradicts (\ref{eq:n-2}). 
Thus, we have \[x_{i_1} + x_{i_2} + \cdots + x_{i_{n-3}} \leq (n-1) + (n-2) + \cdots + 3.\] 
Lastly, consider if $x_1 + x_2 + \cdots + x_{n-1} > (n-1) + \cdots +1$. 
Then \[x_1 + x_2 + \cdots + x_{n-1} + x_n >  n+ (n-1) + \cdots +1,\] which contradicts (\ref{eq:n}). 
Thus, \[x_1 + x_2 + \cdots + x_{n-1} \leq (n-1) + \cdots +1.\]
Collecting these results gives the following inequality description.
\begin{align*}
    1 \leq x_i &\leq n-1, &\text{ for } 1 \leq i \leq n-1,\\
    x_i+x_j &\leq (n-1) + (n-2), &\text{ for } i < j < n,\\
    x_i+x_j+x_k &\leq (n-1) + (n-2) + (n-3), &\text{ for } i < j < k < n,\\
    &\vdots\\
    x_{i_1}+ x_{i_2}+ \cdots + x_{i_{n-3}} &\leq (n-1) + (n-2) + \cdots + 3, &\text{ for } i_1 < i_2 < \cdots < i_{n-2} < n,\\
    x_{1}+ x_{2}+ \cdots + x_{n-1} &\leq  (n-1) + \cdots + 1.
\end{align*}
This is exactly the inequality description for $\PF_{n-1}$. 
Hence, $\PF_{n-1}$ is a facet of $\PF_n$.
Now, as the choice of $i$ was arbitrary, there are $n$ choices for $i$, so there are $n$ copies of $\PF_{n-1}$ appearing as facets of $\PF_n$. 

Now we will show that these are the only occurrences. 
We know that $\PF_{n-1}$ has $(n-1)!\sum\limits_{m=1}^{n-1} \frac{1}{m!}$ many  vertices.
Similar to the proof of Proposition \ref{prop:permutahedron}, we can see that for a hyperplane with $k$ variables, there are $k!(n-k)!\sum\limits_{m=1}^{n-k} \frac{1}{m!}$ many vertices that satisfy it. 
If $k$ is not equal to $1$ or $n$, we have $n\leq \binom{n}{k}$, thus $k!(n-k)!\leq (n-1)!$.
Then since $$k!(n-k)!\sum\limits_{m=1}^{n-k} \frac{1}{m!} < k!(n-k)!\sum\limits_{m=1}^{n-1} \frac{1}{m!} \leq (n-1)! \sum\limits_{m=1}^{n-1} \frac{1}{m!},$$  these hyperplanes cannot correspond to a facet with $(n-1)!\sum\limits_{m=1}^{n-1} \frac{1}{m!}$ many vertices.
If $k =n$, then the vertices that satisfy the equation are exactly all permutations of $(1,2,\dots, n)$, which is $n! > (n-1)! \sum\limits_{m=1}^{n-1} \frac{1}{m!}$ many vertices, hence this hyperplane cannot correspond to $\PF_{n-1}$. 

For a hyperplane of the form $x_i = 1$, there are $\sum_{m=0}^{n-1}\frac{(n-1)!}{m!} \neq \sum_{m=1}^{n-1} \frac{(n-1)!}{m!}$ many vertices that satisfy the equation, hence it also cannot correspond to $\PF_{n-1}$. 
Thus, the only hyperplanes that support a facet of the form $\PF_{n-1}$ are the $n$ mentioned above.
\end{proof}

\end{document}
